\documentclass[11pt]{article}

% ===================== PACKAGES =====================
\usepackage[a4paper,margin=1in]{geometry}
\usepackage{amsmath,amssymb}
\usepackage{enumitem}
\usepackage{fancyhdr}
\usepackage{xcolor}

% ===================== HEADER & FOOTER =====================
\pagestyle{fancy}
\fancyhf{}
\lhead{CSE 400: Fundamentals of Probability in Computing}
\rhead{Lecture 6 Scribe}
\cfoot{\thepage}

% ===================== CUSTOM COMMANDS =====================
\newcommand{\solution}{
    \vspace{0.5cm}
    \noindent\textbf{\textcolor{blue}{Solution:}}
    \vspace{0.2cm}
}

% ===================== TITLE =====================
\title{
    \normalsize School of Engineering and Applied Science (SEAS), Ahmedabad University \\
    \vspace{0.2cm}
    \textbf{CSE 400: Fundamentals of Probability in Computing}\\
    \Large Lecture scribe submission
}
\author{}
\date{}

\begin{document}
\maketitle

\vspace{-2cm}
\begin{center}
    \begin{tabular}{ll}
        \textbf{Name:} & Pranav Jain \\[1.5ex]
        \textbf{Enrollment No:} & AU2420166 \\[1.5ex]
        \textbf{Group:} & s1\_g1\_fin \\[1.5ex]
        \textbf{Email:} & pranav.j@ahduni.edu.in \\[1.5ex]
        \textbf{Lecture:} & 6 \\[1.5ex]
        \textbf{Date of Submission:} & February 7\textsuperscript{th}, 2026 (11:59 PM)
    \end{tabular}
\end{center}

\bigskip

\section{Scope (as on slides)}
Items explicitly covered in the lecture slides used as source:
\begin{enumerate}[left=0pt]
    \item Types of discrete random variables (Bernoulli, Binomial, Geometric, Poisson) — listed.
    \item Probability Mass Function (PMF) — definition and properties.
    \item Example: tossing 3 fair coins (PMF example).
    \item Example: PMF of the form \(p(i)=c\frac{\lambda^i}{i!},\; i=0,1,2,\dots\) — find normalization constant and probabilities.
    \item Bayes' Theorem — statement and worked auditorium example (row/seat selection).
    \item Outline mentions: CDF/PDF, Expectation, moments (variance, etc.) — listed in slides; derivations for these are not present on the provided lecture pages.
\end{enumerate}

\section{Definitions and Notation (as introduced on slides)}
\begin{enumerate}[left=0pt]
    \item \textbf{Random Variable (RV):} variable assigning real values to outcomes of an experiment (illustrated on slides).
    \item \textbf{Discrete random variable:}
    \begin{itemize}
        \item Range (support) \(R_X=\{x_1,x_2,\dots\}\) (finite or countably infinite).
        \item Notation: \(P_X(x_k)=P(X=x_k)\) for \(k=1,2,\dots\).
    \end{itemize}
    \item \textbf{Probability Mass Function (PMF):} For discrete \(X\) with range \(R_X=\{x_k\}\),
    \[
        p_X(x_k)=P(X=x_k),\qquad k=1,2,3,\dots
    \]
    \item \textbf{PMF normalization property:}
    \[
        \sum_k p_X(x_k)=1.
    \]
    \item \textbf{Conditional probability (recap on slides):}
    \[
        P(A\mid B)=\frac{P(A\cap B)}{P(B)}.
    \]
    \item \textbf{Bayes formula (boxed on slides):}
    \[
        P(B_i\mid A)=\frac{P(A\mid B_i)\,P(B_i)}{\sum_j P(A\mid B_j)\,P(B_j)}.
    \]
\end{enumerate}

\section{Worked Example 1: Tossing 3 fair coins (PMF example)}
\textbf{Set-up (from slides):} Toss 3 fair coins. Let \(Y\) denote the number of heads. Range: \(\{0,1,2,3\}\).

\solution

\noindent PMF values (as given on slides):
\begin{align*}
    P(Y=0) &= P(t,t,t) = \frac{1}{8},\\[4pt]
    P(Y=1) &= P(\{t,t,h\},\{t,h,t\},\{h,t,t\}) = \frac{3}{8},\\[4pt]
    P(Y=2) &= P(\{h,h,t\},\{h,t,h\},\{t,h,h\}) = \frac{3}{8},\\[4pt]
    P(Y=3) &= P(h,h,h) = \frac{1}{8}.
\end{align*}

\noindent Normalization check (as on slide):
\[
    1 \;=\; \sum_{i=0}^3 P(Y=i).
\]

\section{Worked Example 2: PMF of the form \(p(i)=c\frac{\lambda^i}{i!}\), \(i=0,1,2,\dots\) (Poisson form)}
\textbf{Given (slide):} \(p(i)=c\frac{\lambda^i}{i!}\) for \(i=0,1,2,\dots\), \(\lambda>0\). Find \(P(X=0)\) and \(P(X>2)\).

\solution

\noindent Use normalization \(\sum_{i=0}^{\infty} p(i)=1\):
\[
    \sum_{i=0}^{\infty} c\frac{\lambda^i}{i!}=1.
\]
Recognize the exponential series (as used on slides):
\[
    \sum_{i=0}^{\infty}\frac{\lambda^i}{i!}=e^{\lambda}.
\]
Thus
\[
    c\,e^{\lambda}=1 \quad\Rightarrow\quad c=e^{-\lambda}.
\]
Hence
\[
    P(X=0)=c\frac{\lambda^0}{0!}=e^{-\lambda}.
\]
For \(P(X>2)\) (complementary sum used on slides):
\[
    \begin{aligned}
        P(X>2) &= 1-P(X\le 2) \\
               &= 1-\big(P(X=0)+P(X=1)+P(X=2)\big) \\
               &= 1 - e^{-\lambda}\Big(1+\lambda+\tfrac{\lambda^2}{2}\Big).
    \end{aligned}
\]

\section{Bayes' Theorem (recap) and Worked Example: Auditorium / Door prize (slides)}
\subsection*{Bayes' Theorem (as boxed on slides)}
\[
    P(B_i\mid A)=\frac{P(A\mid B_i)\,P(B_i)}{\sum_j P(A\mid B_j)\,P(B_j)}.
\]

\subsection*{Auditorium example (as on slides)}
\textbf{Problem statement (from slide):}
\begin{itemize}
    \item Auditorium has 30 rows. Row 1 has 11 seats, Row 2 has 12 seats, Row 3 has 13 seats, \dots, up to Row 30 which has 40 seats (seat counts increase as listed on slide).
    \item Door prize selection: first randomly select a row (each of the 30 rows equally likely), then randomly select a seat within that row (each seat in chosen row equally likely).
\end{itemize}

\noindent Events defined on slides:
\begin{itemize}
    \item \(S_{15}\): seat 15 is selected.
    \item \(R_k\): row \(k\) is selected.
\end{itemize}

\noindent Questions shown on slide:
\begin{enumerate}
    \item Compute \(P(S_{15})\).
    \item Compute \(P(S_{15}\mid R_{20})\).
    \item Compute \(P(R_{20}\mid S_{15})\).
\end{enumerate}

\solution

\noindent (1) Total probability for \(S_{15}\) (as on slides):
\[
    P(S_{15}) \;=\; \sum_{k=1}^{30} P(S_{15}\mid R_k)\,P(R_k).
\]
Here \(P(R_k)=\dfrac{1}{30}\) for each row. \(P(S_{15}\mid R_k)=\dfrac{1}{\text{(number of seats in Row }k\text{)}}\) if seat 15 exists in row \(k\), otherwise 0. The slides evaluate this sum numerically and report:
\[
    P(S_{15}) = 0.0342 \quad\text{(value shown on slide).}
\]

\noindent (2) Conditional probability given row 20 (as on slide):
\[
    P(S_{15}\mid R_{20}) = \frac{1}{\text{number of seats in Row 20}}.
\]

\noindent (3) Posterior via Bayes (as on slide):
\[
    P(R_{20}\mid S_{15})
      \;=\; \frac{P(S_{15}\mid R_{20})\,P(R_{20})}{P(S_{15})}.
\]
Substitute \(P(R_{20})=\tfrac{1}{30}\), \(P(S_{15})=0.0342\) and \(P(S_{15}\mid R_{20})=\dfrac{1}{\text{seats in Row 20}}\) as shown on the slide to obtain the numeric result presented on the slide.

\section{Other items shown on slides (no derivation on provided pages)}
\begin{itemize}[left=0pt]
    \item Discrete vs continuous variables — PMF vs PDF distinction (figure and bullets appear on slides).
    \item CDF / PDF are listed in lecture outline.
    \item Expectation, linearity, moments (variance, skewness, kurtosis) are listed in lecture outline; explicit step-by-step derivations for expectation and moments are not present on the provided Lecture 6 pages.
\end{itemize}

\section*{Boxed formula / Proposition (as on slides)}
\[
    \boxed{ \;P(B_i\mid A)=\dfrac{P(A\mid B_i)\,P(B_i)}{\sum_j P(A\mid B_j)\,P(B_j)}\; }
\]

\bigskip

\noindent\textbf{Note:} This scribe contains only definitions, notation, stated assumptions, theorem/result statements, and the step-by-step proofs/derivations or worked example steps that appear in the provided Lecture 6 slides. No additional explanations, intuition, or new examples have been introduced.

\end{document}
